% ===================================================================
%  పట్టిక సంఖ్య 2 — మానవ శరీరం. --- విరామ సమయం. ఆటలు.
%
%  Traduction, faite en 2026, en telougou d'aujourd'hui, sur l'ido de
%  texto/io. Regles de la colonne : voir 10-pattika-01.tex.
%
%  LE CORPS EST ENTIEREMENT TELOUGOU, ET C'EST LA MEILLEURE PREUVE
%  QUE LA COLONNE N'EST PAS UN CALQUE. Une traduction pressee
%  emprunterait ici au sanskrit, qui a un mot savant pour chaque
%  organe et que le telougou ecrit sans peine. Or le telougou de tous
%  les jours a les siens, dravidiens, et ce sont ceux-la :
%
%    తల       la tete        మొండెం    le tronc
%    పుర్రె    le crane       మెదడు     le cerveau
%    నుదురు   le front       కణతలు     les tempes
%    కన్ను    l'œil          కనుబొమ    le sourcil
%    ముక్కు   le nez         చెంపలు    les joues
%    చెవి     l'oreille      నోరు      la bouche
%    పెదవులు  les levres     నాలుక     la langue
%    దవడ      la machoire    పళ్ళు     les dents
%    మెడ      le cou         గొంతు     la gorge
%    ఛాతీ     la poitrine    గుండె     le cœur
%    ఊపిరితిత్తులు les poumons  పక్కటెముకలు les cotes
%    వీపు     le dos         భుజం      l'epaule
%    పొట్ట    le ventre      కడుపు     l'estomac
%    ప్రేగులు  les intestins  మోచేయి    le coude
%    ముంజేయి  l'avant-bras   మణికట్టు  le poignet
%    వేళ్ళు   les doigts     గోళ్ళు    les ongles
%    సిర      la veine       రక్తం     le sang
%    తొడ      la cuisse      మోకాలు    le genou
%    పిక్క    le mollet      చీలమండ    la cheville
%    అరికాలు  la plante      మడమ       le talon
%    ఎముకలు   les os         చర్మం     la peau
%
%  Et les cinq sens de meme : చూపు, వినికిడి, వాసన, రుచి, స్పర్శ.
%
%  LES JEUX TOMBENT JUSTE, ET C'EST LA QUATRIEME COLONNE OU ILS
%  TOMBENT JUSTE. దాగుడుమూతలు est exactement le cache-cache,
%  కళ్ళగంతల ఆట exactement le colin-maillard, గోళీలు les billes,
%  బొంగరం la toupie, గాలిపటం le cerf-volant — celui de Sankranti,
%  qu'on lance ici par milliers —, ఊయల la balancoire, ఎలుగుబంటి ఆట
%  le jeu de l'ours, కాలిబంతి le football. Aucun n'a demande de
%  periphrase.
%
%  ET LE JEU DE BARRES MANQUE, COMME IL MANQUAIT EN MARATHE. L'Andhra
%  a కబడ్డీ et చెడుగుడు, qui sont d'autres jeux ; ecrire l'un des
%  deux serait traduire le pays, pas le livret. On decrit : అడ్డుల
%  ఆట. Meme trou, meme remede, deuxieme langue indienne d'affilee.
% ===================================================================

%%K t02-apar-1 apar
\VUsaut{4.0mm}
\VUtitre{15.0pt}{60}{పట్టిక సంఖ్య 2}
\VUsaut{2.0mm}
\VUtitre{11.4pt}{0}{మానవ శరీరం. --- విరామ సమయం.}
\VUtitre{11.4pt}{0}{ఆటలు.}

%%K t02-tit-0 sub
\VUtitre{13.2pt}{0}{మానవ శరీరం.}
\VUsaut{2.4mm}
\VUcentre{10.2pt}{0}{\textit{(ప్రకృతి శాస్త్రంలో ఒక సాధారణ పాఠం.)}}
\VUsaut{3.0mm}

%%K t02-tit-1 sub pos=t02-01-1
\VUpk{10.2pt}{40}{తల.}
\VUsaut{3.0mm}

%%K t02-01-1 p
1. --- మానవ శరీరంలోని ముఖ్యమైన భాగాలు : \VUgras{తల}, \VUgras{మొండెం} మరియు \VUgras{అవయవాలు}.

%%K t02-02-1 p
2. --- తలలో రెండు భాగాలు ఉన్నాయి : \VUgras{పుర్రె}, దానిలో \VUgras{మెదడు} ఉంటుంది, దాన్ని
\VUgras{వెంట్రుకలు} కప్పుతాయి; మరియు \VUgras{ముఖం}.

%%K t02-03-1 p
3. --- మన బొమ్మలో పుర్రె కానీ మెదడు కానీ కనిపించడం లేదు; కానీ ముఖంలోని ముఖ్యమైన భాగాలు మనకు
చాలా స్పష్టంగా కనిపిస్తున్నాయి.

%%K t02-04-1 p
4. --- ఇదిగో ముందుగా \VUgras{నుదురు}, ఇది బలమైన బుద్ధిని, ఎప్పుడూ పనిచేసే ఆలోచనను చూపుతుంది.
\VUgras{కణతలు} చాలా విశాలంగా ఉన్నాయి. \VUgras{కన్ను} చురుకుగా, బాగా తెరుచుకుని ఉంది. దట్టమైన
\VUgras{కనుబొమ} మరియు పొడవైన \VUgras{కనురెప్పల వెంట్రుకలు} దానికి నీడనిస్తాయి.

%%K t02-05-1 p
5. --- ఇంకా ఇదిగో \VUgras{ముక్కు} మరియు \VUgras{ముక్కు రంధ్రాలు}. వాసన చూసేందుకు, ఊపిరి
తీసుకునేందుకు ముక్కు మనకు ఉపయోగపడుతుంది. ముక్కుకు ఇరువైపులా \VUgras{చెంపలు}, ఇంకా దూరంగా
\VUgras{చెవులు} ఉన్నాయి. ముఖం కింది భాగం \VUgras{నోరు} మరియు \VUgras{గడ్డంతో} ఏర్పడుతుంది.

%%K t02-06-1 p
6. --- అవి మనకు అంత బాగా కనిపించడం లేదు, ఎందుకంటే \VUgras{మీసం} మరియు \VUgras{చెంపల గడ్డం}
వాటిని మన నుంచి దాచేస్తున్నాయి. కానీ నోటిలో రెండు \VUgras{పెదవులు}, \VUgras{పైదవడ} మరియు
\VUgras{కిందిదవడ} వాటి \VUgras{పళ్ళతో}, ఇంకా \VUgras{నాలుక} ఉంటాయని మనకు తెలుసు. నోటితో మనం
మాట్లాడతాం, తింటాం. నాలుకతో ఆహారం రుచి చూస్తాం.

%%K t02-07-1 p
7. --- కన్ను, చెవి, ముక్కు మరియు నోరు — వేళ్ళలాగే — ఐదు ఇంద్రియాల అవయవాలు : \VUgras{చూపు},
\VUgras{వినికిడి}, \VUgras{వాసన}, \VUgras{రుచి}, \VUgras{స్పర్శ}.

%%K t02-08-1 p
8. --- కాబట్టి మానవ శరీరంలో అత్యంత ఆసక్తికరమైన భాగాల్లో తల ఒకటి.

%%K t02-tit-2 sub
\VUsaut{4.4mm}
\VUpk{10.2pt}{40}{మొండెం.}
\VUsaut{3.0mm}

%%K t02-09-1 p
9. --- \VUgras{మెడ} తలను మొండెంతో కలుపుతుంది. అందులో \VUgras{గొంతు} మరియు
\VUgras{మెడ వెనుకభాగం} ఉంటాయి.

%%K t02-10-1 p
10. --- మొండెంలో ఎన్నో ముఖ్యమైన అవయవాలు కూడా ఉన్నాయి. \VUgras{ఛాతీలో} \VUgras{ఊపిరితిత్తులు}
ఉన్నాయి — వాటితో మనం ఊపిరి తీసుకుంటాం — మరియు \VUgras{గుండె}, అది జీవానికి కేంద్రం. గుండె
కొట్టుకోవడం ఆగిపోతే జీవి చనిపోతుంది.

%%K t02-11-1 p
11. --- \VUgras{పక్కటెముకలు} ఛాతీని నిలబెడతాయి.

%%K t02-12-1 p
12. --- మొండెంలోని మిగతా భాగాలు \VUgras{ప్రక్క}, \VUgras{వీపు}, \VUgras{భుజాలు} మరియు
\VUgras{పొట్ట} (కడుపుభాగం). నిద్రపోయేందుకు మనం ప్రక్కకో వీపుపైనో పడుకుంటాం, బరువులను భుజంపై
మోస్తాం.

%%K t02-13-1 p
13. --- పొట్టలో \VUgras{కడుపు} మరియు \VUgras{ప్రేగులు} ఉన్నాయి. ఆహారాన్ని జీర్ణం చేసేందుకు అవి
మనకు ఉపయోగపడతాయి.

%%K t02-tit-3 sub
\VUsaut{4.4mm}
\VUpk{10.2pt}{40}{అవయవాలు.}
\VUsaut{3.0mm}

%%K t02-14-1 p
14. --- మనకు నాలుగు \VUgras{అవయవాలు} ఉన్నాయి : రెండు \VUgras{చేతులు} మరియు రెండు
\VUgras{కాళ్ళు}. చేతులతో మనం వస్తువులను తీసుకుంటాం, పట్టుకుంటాం, ఎత్తుతాం, పోగుచేస్తాం;
కాళ్ళతో నిలబడతాం, నడుస్తాం, పరుగెడతాం.

%%K t02-15-1 p
15. --- \VUgras{భుజం} చేతిని శరీరంతో కలుపుతుంది. చాలా బలమైన ఒక \VUgras{కండరం} — కండరపు బలం —
చేతిని కదిలిస్తుంది, ఆ చేతిని \VUgras{మోచేయి} \VUgras{ముంజేతితో} కలుపుతుంది. \VUgras{మణికట్టు}
\VUgras{చేతికి} కీలు అవుతుంది, ఆ చేయి \VUgras{వేళ్ళతో} మరియు \VUgras{గోళ్ళతో} ముగుస్తుంది.
చేతిపై ఒక \VUgras{సిర} (మరియు ధమని) మనకు కనిపిస్తుంది, దానిలో \VUgras{రక్తం} ప్రవహిస్తుంది.

%%K t02-16-1 p
16. --- కాలిలోని భాగాలు : \VUgras{తొడ}, \VUgras{మోకాలు} మరియు \VUgras{మోకాలి వెనుకభాగం},
\VUgras{పిక్క} — దాని \VUgras{మాంసాన్ని} \VUgras{చర్మం} కప్పుతుంది —, \VUgras{చీలమండ} మరియు
\VUgras{పాదం}. కాలి \VUgras{ఎముకలు} చాలా లావుగా, చాలా బలంగా ఉంటాయి. పాదం చివర
\VUgras{కాలివేళ్ళు} ఉన్నాయి. మిగతా భాగాలు \VUgras{అరికాలు} మరియు \VUgras{మడమ}.

%%K t02-17-1 p
17. --- మానవ శరీరానికి దాదాపు పూర్తి వర్ణన ఇది.

%%K t02-17-2 p
\textit{గమనిక.} ---
\textit{« నాకు... (తల మొదలైనవి) నొప్పిగా ఉంది » అనే వాక్యం సాయంతో ఉపాధ్యాయుడు ఈ పదజాలం మొత్తాన్ని మంచి లేదా చెడు}
\VUgras{శారీరక స్థితి} \textit{దృష్టి నుంచి మళ్ళీ వాడుకోగలడు.}

%%K t02-tit-4 sub
\VUsaut{5.0mm}
\VUpk{10.2pt}{40}{వ్యాయామం.}
\VUsaut{2.4mm}
\VUcentre{10.2pt}{0}{\textit{(ఒక విద్యార్థి చెప్పిన కథనం.)}}
\VUsaut{3.0mm}

%%K t02-18-1 p
18. --- మెత్తగా, చురుకుగా, ఆరోగ్యంగా ఉండాలంటే మన కాళ్ళకూ చేతులకూ వ్యాయామం ఇవ్వాలి. అందుకే మనం
ఆడతాం, \VUgras{వ్యాయామం} చేస్తాం.

%%K t02-19-1 p
19. --- వారంలో ఈ రెండు వ్యాయామాలూ ఉన్నత పాఠశాల ఆవరణలో జరుగుతాయి (పట్టిక సంఖ్య 1 చూడండి). కానీ
గురువారం, ఆదివారం మనం పెద్ద పచ్చిక బయలుకు వెళ్తాం, అక్కడ పరుగెత్తేందుకూ ఆడుకునేందుకూ మనకు చోటు
ఉంటుంది.

%%K t02-20-1 p
20. --- మన గురువులూ మన తల్లిదండ్రులూ అప్పుడప్పుడూ మనతో అక్కడికి వస్తారు; కానీ వ్యాయామాలను
నడిపేది \VUgras{వ్యాయామ ఉపాధ్యాయుడే} \textsuperscript{(12)}.

%%K t02-21-1 p
21. --- పచ్చిక బయలులో, ప్రవేశద్వారానికి ఎడమవైపు \VUgras{వ్యాయామ సాధనాలు} \textsuperscript{(1)}
ఉన్నాయి; \VUgras{నిచ్చెనతో} \textsuperscript{(2)} మనం ఎక్కుతాం, \VUgras{ఉంగరాలకూ}
\textsuperscript{(3)} \VUgras{త్రాపెజ్‌కూ} \textsuperscript{(4)} వేలాడి తలకిందులవుతాం,
\VUgras{తాడు నిచ్చెన} \textsuperscript{(5)} లేదా \VUgras{ముడుల తాడు} \textsuperscript{(6)}
చివరిదాకా చేతులతో ఎగబాకుతాం.

%%K t02-22-1 p
22. --- ఈలోగా మన స్నేహితులు \VUgras{చెక్క గుర్రాన్ని} \textsuperscript{(7)} లేదా
\VUgras{సమాంతర కమ్మీలను} \textsuperscript{(11)} దూకుతారు. మరికొందరు \VUgras{బాక్సింగ్}
\textsuperscript{(8)} చేస్తారు లేదా ముఖకవచంతో, \VUgras{ఫాయిల్‌తో} \VUgras{కత్తియుద్ధం}
\textsuperscript{(9)} చేస్తారు. చివరగా, మరికొందరు \VUgras{బరువులు ఎత్తుతారు}
\textsuperscript{(10)}.

%%K t02-tit-5 sub
\VUsaut{4.6mm}
\VUpk{10.2pt}{40}{ఆటలు.}
\VUsaut{3.0mm}

%%K t02-23-1 p
23. --- వ్యాయామం చేసేవాళ్ళ చుట్టూ అబ్బాయిలూ అమ్మాయిలూ రకరకాల \VUgras{ఆటలు} ఆడతారు. అమ్మాయిలు
తమ \VUgras{చక్రంతో} \textsuperscript{(14)} ఆడుకుంటారు; \VUgras{తాడుతో} \textsuperscript{(13)}
దూకుతారు, లేదా తమ తమ్ముళ్ళనూ చిన్నాన్న కొడుకులనూ \VUgras{క్రొకె} \textsuperscript{(15)} ఆటకు
పిలుస్తారు. ప్రత్యర్థి తన \VUgras{బంతిని} సులువుగా వంపు గుండా దాటిస్తానని అనుకున్న క్షణంలోనే
దాన్ని తమ చెక్క సుత్తితో దూరంగా కొట్టడంలో వాళ్ళకు మజా.

%%K t02-24-1 p
24. --- అబ్బాయిలకు కండల ఆటలే ఇష్టం. \VUgras{స్కిటిల్స్} \textsuperscript{(16)} ఆడటం వాళ్ళకు
ఇష్టం, వాటిని \VUgras{బంతితో} \textsuperscript{(17)} నేర్పుగా పడగొడతారు. \VUgras{బొంగరం}
\textsuperscript{(18)}, \VUgras{బిల్‌బోకె} \textsuperscript{(19)}, చివరకు \VUgras{కప్ప ఆట}
\textsuperscript{(20)} కూడా వాళ్ళు తక్కువ చేయరు. కానీ వాళ్ళకు అన్నిటికంటే ఇష్టమైనవి
\VUgras{గోళీలు} \textsuperscript{(21)} మరియు \VUgras{గోడ బంతి} \textsuperscript{(22)},
ఎందుకంటే \VUgras{హరితగృహం} \textsuperscript{(26)} గోడ తేలికైన బంతిని తిప్పికొట్టేందుకు చాలా
బాగుంటుంది.

%%K t02-25-1 p
25. --- తన \VUgras{కర్రకాళ్ళపై} \textsuperscript{(24)} నడిచే చిన్న జాన్ చాలా నేర్పరి,
సమతూకాన్ని బాగా నిలబెట్టుకోగలడు. అతని దగ్గరే కొందరు స్నేహితులు ఆ హరితగృహంలోనూ \VUgras{కంచె}
వెనుకా \VUgras{దాగుడుమూతలు} \textsuperscript{(25)} ఆడుతున్నారు. మరికొందరికి
\VUgras{కొట్టినవాడిని కనిపెట్టే ఆట} \textsuperscript{(28)} లేదా \VUgras{ఎగిరే బంతి}
\textsuperscript{(29)} ఇష్టం — అది ఒక \VUgras{రాకెట్} నుంచి ఇంకో రాకెట్‌కు ఎగురుతూ ఉంటుంది.

%%K t02-noto-1 noto
\VUnotes{143.2mm}{%
(*) పిల్లల ఆటలకు తరచూ పూర్తిగా జాతీయమైన పేర్లు ఉంటాయి. వాటిని ఇడోలో వీలైనంత చక్కగా
పేర్కొనేందుకు మేము ప్రయత్నించాం — ఒకోసారి ఆ ఆటను గుర్తించే చర్య నుంచే ప్రేరణ తీసుకుని, ఒకోసారి
ఆ దేశపు పేరు మరీ అన్యాయంగా లేనప్పుడు దాన్నే ఎంచుకుని. ఉదాహరణకు \VUgras{పీపా ఆట} (జర్మన్,
ఫ్రెంచ్) అంటే \VUgras{కప్ప ఆట} \textsuperscript{(20)} (ఇంగ్లిష్), అందులో ఆడేవాళ్ళు తమ గుండ్రటి
చక్రాలను, రంధ్రాలున్న పెట్టె (పీపా) పై నోరు తెరుచుకుని నిలబడ్డ లోహపు కప్ప నోటిలోకి విసరడానికి
ప్రయత్నిస్తారు. \VUgras{భుజాల దూకుడు} \textsuperscript{(30)} \VUgras{వీపు దూకుడుకు} ఒక్క
విషయంలోనే వేరు : తలపైనుంచి దూకేందుకు ఆడేవాళ్ళు తమ చేతులను జోడీవాడి భుజాలపై ఆనిస్తారు, «
\VUgras{వీపు దూకుడులో} » మాత్రం చేతులను అతని వీపుపైనే ఆనిస్తారు. --- లేదా పాల్గొనేవాళ్ళలో
కొందరు వంగి ఉంటారు, మిగతావాళ్ళు చేతులను భుజంపై ఆనించి వాళ్ళ వీపుపైనుంచి దూకుతారు; మొదటివాళ్ళు
వంగిపోకుండా ఆ దెబ్బను భరించాలి, రెండోవాళ్ళు సమతూకం కోల్పోకుండా దూకాలి (పట్టికనే చూడండి).%
}

%%K t02-26-1 p
26. --- పచ్చిక బయలు మధ్యలో రెండు చెట్లు ఉన్నాయి. వాటిలో ఒకదాని మొదట్లో ఏడుగురు లేదా ఎనిమిది
మంది అబ్బాయిలు \VUgras{భుజాల దూకుడు} \textsuperscript{(30)} (*) ఆట మొదలుపెట్టారు. ఈ ఆట అన్ని
దేశాల్లోనూ తెలియదు; కానీ \VUgras{వీపు దూకుడు} ఏమిటో అందరికీ తెలుసు — ఈసారి పిల్లలు అది ఆడటం
లేదు.

%%K t02-tit-6 sub
\VUsaut{5.0mm}
\VUpk{10.2pt}{40}{బయలు ఆటలు.}
\VUsaut{3.4mm}

%%K t02-27-1 p
27. \VUgras{కళ్ళగంతల ఆట} \textsuperscript{(31)} కూడా చాలా సరదాగా ఉంటుంది; అమ్మాయిలూ అబ్బాయిలూ
ఒకరి తర్వాత ఒకరు కళ్ళకు \VUgras{కళ్ళగంత} కట్టుకుని, పట్టుబడేలా తిరిగే నేర్పులేనివాళ్ళను
పట్టుకుంటారు.

%%K t02-28-1 p
28. --- పచ్చిక బయలు మధ్యలో ఇదిగో చాలా ఫ్రెంచ్ అయిన, కొంచెం ఉద్ధతమైన ఆట : అదే
\VUgras{ఎలుగుబంటి ఆట} \textsuperscript{(32)}, అది \VUgras{గాలిపటం} \textsuperscript{(33)} అనే
అంత ప్రశాంతమైన ఆటకంటే, చివరకు \VUgras{టెన్నిస్} \textsuperscript{(34)} అనే ఇంగ్లిష్ ఆటకంటే
కూడా ఎంతో ఎక్కువ గోల చేస్తుంది.

%%K t02-29-1 p
29. --- ఈ చివరి ఆటే పెద్ద విద్యార్థుల ఆనందం. మన ఉపాధ్యాయులే దాన్ని ఇష్టంగా ఆడతారు. దానికి
గొప్ప నేర్పూ చురుకుదనమూ కావాలి, అది నాకు చాలా నచ్చుతుంది. \VUgras{ఊయల} \textsuperscript{(35)}
ఆటలను నేను అమ్మాయిలకు వదిలేస్తాను.

%%K t02-30-1 p
30. --- ఇంకో ఇంగ్లిష్ ఆట నాకు నచ్చదు, అది ఒకవేళ ప్రమాదకరం కావచ్చు.

%%K t02-30-1b p
నేను \VUgras{కాలిబంతి} \textsuperscript{(36)} గురించే చెప్పాలనుకుంటున్నాను, అది మా అన్నయ్యకు
ఆనందం కలిగిస్తుంది. అంతే ఆరోగ్యకరమైన, నిజంగా తక్కువ కర్కశమైన మన పాత \VUgras{అడ్డుల ఆటే}
\textsuperscript{(38)} నాకు ఇష్టం.

%%K t02-tit-7 sub
\VUsaut{4.6mm}
\VUpk{10.2pt}{40}{సైకిల్ తొక్కడం మరియు బంతి ఆట.}
\VUsaut{3.0mm}

%%K t02-31-1 p
31. --- \VUgras{సైకిల్‌పై} \textsuperscript{(37)} పచ్చిక బయలు చుట్టూ తిరగడమూ నాకు ఇష్టమే.

%%K t02-32-1 p
32. --- వచ్చే సంవత్సరం నేను \VUgras{గుర్రపు స్వారీ} \textsuperscript{(39)} నేర్చుకుంటాను.
హుందాగా అడుగులు వేసే లేదా దౌడు తీసే, అడ్డంకులను దూకి దాటే గుర్రం ఎంత అందంగా ఉంటుంది!

%%K t02-33-1 p
33. --- వేసవిలో మనం \VUgras{ఈత} \textsuperscript{(40)} నేర్చుకుంటాం. నదిలోని స్వచ్ఛమైన చల్లని
నీళ్ళలో మనం ఆనందంగా మునుగుతాం, స్నానం చేస్తాం. అప్పుడప్పుడూ చివరగా ఒక కాగితపు \VUgras{బెలూన్}
\textsuperscript{(41)} ఎగరేస్తాం, వేడిగాలితో నిండగానే అది గంభీరంగా గాలిలోకి లేస్తుంది.

%%K t02-34-1 p
34. --- ఇంకా ఎన్నో ఆటలు ఉన్నాయి, కానీ వాటిని లెక్కపెట్టడం నాకు ముగియదు; ఈ సంక్షిప్త వర్ణనను
బట్టి, గురువారాల్లో మన అందమైన పచ్చిక బయలులో ఎవరికీ పెద్దగా విసుగు రాదని అర్థమవుతుంది.

%%K t02-34-2 p
గమనిక. ---
\textit{అత్యంత ముఖ్యమైన ఆటల్లో ప్రతిదాన్నీ మరింత వివరంగా వర్ణిస్తూ ఉపాధ్యాయుడు ఈ అధ్యాయాన్ని సులువుగా పూర్తి చేయగలడు.}
